%% BigDataSmallPrice – Abschlussbericht
%% ZHAW Semesterprojekt Big Data (DS23t), FS 2026
%% Bachmann Ryan · Dinis Silva Miguel · Ruchti Gian

\documentclass[sigconf]{acmart}

\usepackage[ngerman]{babel}
\usepackage[T1]{fontenc}
\usepackage{booktabs}
\usepackage{multirow}
\usepackage{tikz}
\usetikzlibrary{arrows.meta, positioning}
\usepackage{pgfplots}
\pgfplotsset{compat=1.18}
\usepackage{microtype}

\setcopyright{none}
\acmDOI{}
\acmISBN{}
\acmPrice{}
\acmConference[DS23t]{ZHAW Semesterprojekt -- Big Data}{FS~2026}{Winterthur, Schweiz}
\acmYear{2026}
\copyrightyear{2026}

%% ============================================================
\title{BigDataSmallPrice: Eine Big-Data-Pipeline zur Vorhersage
       dynamischer Stromtarife in Winterthur}

\author{Ryan Bachmann}
\affiliation{%
  \institution{ZHAW School of Engineering}
  \city{Winterthur}
  \country{Schweiz}}
\email{bachmrya@students.zhaw.ch}

\author{Miguel Dinis Silva}
\affiliation{%
  \institution{ZHAW School of Engineering}
  \city{Winterthur}
  \country{Schweiz}}
\email{dinismig@students.zhaw.ch}

\author{Gian Ruchti}
\affiliation{%
  \institution{ZHAW School of Engineering}
  \city{Winterthur}
  \country{Schweiz}}
\email{ruchtgia@students.zhaw.ch}

%% ============================================================
\begin{abstract}
Das revidierte Schweizer Stromversorgungsgesetz (StromVG, Art.~8b) verpflichtet alle
Netzbetreiber ab 2026, ihren Kunden auf Wunsch dynamische Stromtarife anzubieten.
Dynamische Tarife reagieren unmittelbar auf Grosshandelspreise an der EPEX~SPOT und auf
die momentane Auslastung des lokalen Verteilnetzes -- für Endverbraucher ohne geeignete
Prognosewerkzeuge eine kaum planbare Situation.
Wir präsentieren \textsc{BigDataSmallPrice}, eine vollautomatische End-to-End-Big-Data-Pipeline,
die täglich aktuelle 24-Stunden-Vorhersagen für dynamische Netz- und Energietarife
am Standort Winterthur berechnet und über ein interaktives Web-Dashboard bereitstellt.
Das System integriert täglich Daten aus fünf öffentlichen Quellen mit heterogenen Formaten
und Latenzen in einer TimescaleDB-Zeitreihendatenbank mit über fünf Millionen
Datenpunkten. Apache-Airflow-DAGs orchestrieren ETL, Feature-Engineering sowie
Modelltraining vollständig automatisiert. Zwei XGBoost-Modelle prognostizieren
die Winterthurer Netzlast (Modell~A) und den EPEX-CH-Day-Ahead-Spotpreis (Modell~B);
eine regelbasierte Formel-Engine berechnet daraus kundenverständliche Tarife in
Rappen pro Kilowattstunde. Modell~A erreicht einen MAPE von 1{,}51\,\%, Modell~B von
10{,}6\,\% -- beide weit unterhalb der definierten Qualitätsgrenzen von 8\,\%
bzw.\ 15\,\%. Die REST-API liefert eine 24-Stunden-Prognose in unter 100\,ms.
\end{abstract}

\keywords{Dynamische Stromtarife, Zeitreihenprognose, XGBoost, Apache Airflow,
          TimescaleDB, Big Data, EPEX, Netzlastvorhersage, Schweiz}

\maketitle

%% ============================================================
\section{Einleitung}

\subsection{Kontext: Energiewende und dynamische Preise}

Die Energiewende verändert die Struktur des europäischen Strommarkts grundlegend.
Der rasch wachsende Anteil fluktuierender erneuerbarer Energien -- insbesondere Windkraft
und Photovoltaik -- erzeugt starke, schwer vorhersagbare Schwankungen in Angebot und
Nachfrage. Diese Schwankungen schlagen sich unmittelbar in den Grosshandelspreisen an der
europäischen Strombörse EPEX~SPOT nieder~\cite{weron2014}: Während an windstarken Mittagen
negative Preise entstehen können, erreichen Spotpreise bei winterlichen Dunkelflauten
mehrere hundert Euro pro Megawattstunde. Die Schweiz ist über vier Kupplungspunkte
(Deutschland, Frankreich, Österreich, Italien) eng in diesen volatilen Markt eingebunden;
der Schweizer EPEX-CH-Preis folgt dem deutschen Markt in normalen Zeiten nahezu
identisch, wird aber durch alpine Wasserkraftkapazitäten und Netzengpässe regional
entkoppelt.

Gleichzeitig entstehen auf der Verbrauchsseite neue steuerbare Lasten: Elektrofahrzeuge,
Wärmepumpen und Heimspeicher bieten das Potenzial zur zeitlichen Lastverschiebung.
Dynamische Stromtarife sind das marktbasierte Steuerinstrument, das Preissignale
direkt an Endkunden weitergibt und so wirtschaftliche Anreize zur Lastflexibilität schafft.

\subsection{Regulatorischer Rahmen: StromVG Art.~8b}

Das revidierte Bundesgesetz über die Stromversorgung (StromVG) verpflichtet mit
dem neuen Artikel~8b ab dem 1.~Januar~2026 alle Schweizer Netzbetreiber,
ihren Kunden auf Anfrage einen dynamischen Tarif anzubieten~\cite{stomvg2023}.
Der Netznutzungstarif soll dabei die tatsächliche zeitliche Auslastung des
Verteilnetzes widerspiegeln, der Energiepreis die Beschaffungskosten am Grosshandelsmarkt.
Erste Pilotprodukte existieren bereits: EKZ (Zürcher Unterland) und CKW
(Zentralschweiz) haben 2025 dynamische Tarife eingeführt~\cite{ekz2025}.
Diese bieten Kunden, die ihren Verbrauch in Niedrigtarifzeiten verschieben können,
messbare Einsparpotenziale -- typischerweise 15--30\,\% der Energiekosten.

\subsection{Das Problem: Fehlende Planbarkeit für Endverbraucher}

Trotz des regulatorischen Fortschritts bestehen für Endverbraucher zwei fundamentale
Barrieren: Erstens sind die Tarife hochvolatil und nur bis 24~Stunden im Voraus bekannt
(EPEX Day-Ahead-Preise werden täglich um ca.\ 13:00~Uhr für den Folgetag veröffentlicht).
Zweitens fehlen öffentlich zugängliche, benutzerfreundliche Prognosewerkzeuge vollständig.
Die bestehenden Netzbetreiber-APIs~\cite{ekz2025} liefern aktuelle Tarife, aber keine
Vorhersagen. Ein Haushalt mit Wärmepumpe oder Elektroauto kann damit keine fundierte
Entscheidung treffen, wann das Gerät betrieben werden soll.

Die wissenschaftliche Literatur zur Kurzfristprognose von Grosshandelspreisen~\cite{weron2014,lago2018}
und zur probabilistischen Lastvorhersage~\cite{hong2016} liefert zwar methodische Grundlagen,
fokussiert aber auf den Übertragungsnetzbetrieb und nicht auf die kundenseitige
Tarifberechnung. Ein integriertes System, das Prognose und Tarifberechnung
durchgängig verbindet und Endkunden zugänglich macht, fehlte bislang.

\subsection{Ansatz und Beiträge}

Wir entwickeln \textsc{BigDataSmallPrice} als Antwort auf diese Lücke: eine
vollautomatisierte, containerisierte Big-Data-Pipeline, die täglich aktuelle
24\,h-Tarifprognosen berechnet und über ein Web-Dashboard zugänglich macht.
Die Hauptbeiträge dieser Arbeit sind:

\begin{itemize}
  \item \textbf{Datenintegration:} Ein ETL-System, das täglich fünf heterogene
    öffentliche Quellen (XML, JSON, CSV, proprietäre APIs) in einer skalierbaren
    Zeitreihendatenbank konsolidiert und über 5~Millionen Datenpunkte verwaltet.
  \item \textbf{ML-Modelle:} Zwei XGBoost-Modelle für Netzlast- und EPEX-Preisvorhersage,
    kalibriert mit Winterthurer Lokaldaten und deutschen Wettermerkmalen als
    stärksten externen Preistreibern.
  \item \textbf{Tarifberechnung:} Eine transparente, parametrisierbare Formel-Engine,
    die Prognoseausgaben nach dem Vorbild lokaler Netzbetreiber in Rp./kWh-Tarife
    überführt.
  \item \textbf{Open-Source-Pipeline:} Eine vollständige, reproduzierbare
    Produktions\-pipeline auf Basis ausschliesslich quelloffener Komponenten
    (Airflow, TimescaleDB, FastAPI, XGBoost), demonstriert als Semesterprojekt
    durch ein dreiköpfiges Studierendenteam.
\end{itemize}

%% ============================================================
\section{Hintergrund und Problemstellung}

\subsection{Struktur des Schweizer Strommarkts}

Der Schweizer Strommarkt gliedert sich in Erzeugung, Übertragungsnetz (Swissgrid),
Verteilnetze (lokale Betreiber wie EKZ, Stadtwerk Winterthur) und Endkunden.
Grossverbraucher kaufen Strom direkt am EPEX-Spotmarkt; Kleinkunden sind
regulierten Tarifen ihrer lokalen Netzbetreiber ausgesetzt.
Ein dynamischer Tarif setzt sich aus drei Komponenten zusammen:
\textit{Netznutzungsentgelt} (zeitvariabel, abhängig von der Verteilnetzlast),
\textit{Energiepreis} (zeitvariabel, abgeleitet vom EPEX-Day-Ahead-Preis) sowie
fixen staatlichen Abgaben (KEV-Abgabe, Netzkosten Swissgrid, Gemeindegebühr),
die sich nicht ändern. In diesem Projekt modellieren wir die beiden variablen
Komponenten.

\subsection{Tarifformeln und Regulierungsvorgaben}

Aus der Analyse der EKZ-, CKW- und Groupe-E-APIs leiten wir die folgenden
Berechnungsformeln ab. Der Netzpreis folgt einer quadratischen Funktion der
normierten Netzlast:
\begin{equation}
  P_{\mathrm{Netz}} = \mathrm{clip}\!\left(\alpha \cdot \tilde{L}^{2},\;
    \tau_{\mathrm{N}} - 5,\; \tau_{\mathrm{N}} + 15\right)
  \label{eq:netzpreis}
\end{equation}
\begin{equation}
  \tilde{L} = \mathrm{clip}\!\left(\frac{L - L_{\min}}{L_{\max} - L_{\min}},\; 0,\; 1\right)
  \label{eq:lnorm}
\end{equation}
Die quadratische Form stärker gewichtet Extremlastsituationen und schafft damit
einen überproportionalen Anreiz zur Lastverschiebung bei Netzspitzen -- dies entspricht
dem im Verteilnetzbetrieb etablierten Prinzip der Grenzkostenbepreisung.
Die Clipping-Grenzen halten den Tarif in einem regulatorisch vertretbaren Band
um den Standardtarif $\tau_{\mathrm{N}} = 10$\,Rp./kWh.
Mit $L_{\min} = 150$\,kWh (Sommer-Nacht-Minimum) und $L_{\max} = 900$\,kWh
(Winter-Spitzenlast Winterthur) und $\alpha = 15$\,Rp./kWh ergibt sich die in
Abbildung~\ref{fig:tariff_curve} gezeigte Kurve.

\begin{figure}[t]
  \centering
  \begin{tikzpicture}
    \begin{axis}[
      width=\columnwidth,
      height=4.0cm,
      xlabel={Normierte Netzlast $\tilde{L}$},
      ylabel={Netzpreis [Rp./kWh]},
      xmin=0, xmax=1, ymin=0, ymax=17,
      xtick={0,0.25,0.5,0.75,1.0},
      ytick={0,5,10,15},
      grid=major,
      grid style={dotted,gray!50},
      legend pos=north west,
      legend style={font=\scriptsize},
      tick label style={font=\scriptsize},
      label style={font=\small}
    ]
    \addplot[thick, blue!70!black, domain=0:1, samples=200] {15*x^2};
    \addlegendentry{$P_{\mathrm{Netz}} = 15\tilde{L}^{2}$}
    \addplot[dashed, red!70!black, domain=0:1, samples=2] {10};
    \addlegendentry{Standardtarif (10\,Rp./kWh)}
    \addplot[dotted, black!60, domain=0:1, samples=2] {15};
    \addlegendentry{Maximum (15\,Rp./kWh)}
    \end{axis}
  \end{tikzpicture}
  \caption{Netzpreis als quadratische Funktion der normierten Last $\tilde{L}$.
           Bei $\tilde{L} \approx 0{,}82$ wird der Standardtarif von 10\,Rp./kWh
           überschritten; der Peak von 15\,Rp./kWh liegt bei Volllast.}
  \label{fig:tariff_curve}
\end{figure}

Der Energiepreis leitet sich linear aus dem EPEX-Grosshandelspreis ab:
\begin{equation}
  P_{\mathrm{Energie}} = \mathrm{clip}\!\left(k_{pe} \cdot \frac{p_{\mathrm{EPEX}}}{10}
    + k_{le},\; \tau_{\mathrm{E}} - 5,\; \tau_{\mathrm{E}} + 5\right)
  \label{eq:energiepreis}
\end{equation}
mit dem EPEX-Durchreichefaktor $k_{pe} = 0{,}15$, dem fixen Aufschlag
$k_{le} = 2{,}0$\,Rp./kWh, dem Standardenergiepreis $\tau_{\mathrm{E}} = 8$\,Rp./kWh
und dem Umrechnungsfaktor 10 (1~EUR/MWh $\approx$ 1~Rp./kWh bei CHF-EUR-Kurs~1{,}0).
Der Faktor $k_{pe} < 1$ spiegelt wider, dass Netzbetreiber EPEX-Risiken nur
teilweise an Kunden weitergeben. Der Gesamttarif lautet:
\begin{equation}
  P_{\mathrm{gesamt}} = P_{\mathrm{Netz}} + P_{\mathrm{Energie}} \quad [\text{Rp./kWh}]
  \label{eq:gesamttarif}
\end{equation}

\subsection{Technische Grundlagen und Technologiewahl}

\textbf{Zeitreihendatenbank.}
Für die Datenhaltung betrachteten wir drei Alternativen: plain PostgreSQL, InfluxDB
und TimescaleDB. Plain PostgreSQL war wegen fehlender automatischer Zeitreihenpartitionierung
ungeeignet: Tabellen mit mehreren Millionen Zeitreihenpunkten zeigen ohne
manuelle Partitionierung erhebliche Performance-Einbussen bei Zeitbereichsabfragen.
InfluxDB verfügt zwar über hervorragende Time-Series-Performance, erfordert
aber die Flux-Abfragesprache -- was die Integration komplexer SQL-Joins für das
Feature-Engineering erheblich erschwert hätte. \textbf{TimescaleDB}~\cite{freedman2019}
bietet als PostgreSQL-Erweiterung das Beste beider Welten: automatisch partitionierte
Hypertables mit bis zu 20-fach schnelleren Zeitbereichsabfragen bei vollständiger
SQL-Kompatibilität. Dies erlaubt es, die Feature-Engineering-Logik komplett
in SQL-Views zu formulieren, was Konsistenz und Wartbarkeit gegenüber
Python-seitiger Datenbanklogik erheblich verbessert.

\textbf{Gradientverstärkung mit XGBoost.}
Bei der Modellwahl verglichen wir Prophet (Meta), SARIMA, LSTMs und
XGBoost~\cite{chen2016xgboost}. Prophet setzt auf dekomponierbare Trend-Saisonalitäts-Modelle
und ist für externe Kovariablen (Wetter, Feiertage) nur eingeschränkt geeignet.
SARIMA-Modelle setzen Stationarität voraus und können Wettervariablen und
andere multivariate Prädiktoren nicht nativ einbinden. LSTMs erfordern erheblich
mehr Trainingsdaten, sind schwerer zu interpretieren und in der Hyperparametersuche
deutlich aufwändiger. \textbf{XGBoost} hingegen behandelt beliebig gemischte
Feature-Typen (Zeitstempel-Features, Wetterfeatures, Preis-Lags) nativ, braucht keine
Stationarität, unterstützt Early~Stopping und liefert Feature-Importances für die Interpretierbarkeit --
für diesen Anwendungsfall die überzeugendste Kombination aus Genauigkeit,
Trainingseffizienz und Interpretierbarkeit.

\textbf{Workflow-Orchestrierung.}
Gegenüber einfachen Cron-Jobs bietet \textbf{Apache Airflow}~\cite{airflow}
DAG-basierte Aufgabenabhängigkeiten, Task-Level-Logging, eine Web-Oberfläche
für Monitoring und konfigurierbare Retry-Mechanismen. Im Vergleich zu neueren
Alternativen wie Prefect oder Dagster weist Airflow die grösste Community,
umfangreichste Dokumentation und einfachste Deployment-Integration mit bestehenden
PostgreSQL-Umgebungen auf.

\subsection{Formale Problemformulierung}

Gegeben sei eine Zeitreihe historischer Beobachtungen:
$(L_t, p_t, \mathbf{w}_t, \mathbf{x}_t)_{t \leq T}$, wobei $L_t$ die Netzlast
in kWh, $p_t$ der EPEX-CH-Day-Ahead-Preis in EUR/MWh, $\mathbf{w}_t$ ein
Wettervektor und $\mathbf{x}_t$ weitere exogene Merkmale (Kalender, Hydro,
Erzeugungsdaten) zum Zeitpunkt $t$ sind. Gesucht sind zwei Vorhersagefunktionen:
\begin{align}
  \hat{L}_{T+h} &= f_A\!\left(L_{T-k}, \mathbf{w}_{T+h}^{\mathrm{fcast}},
                              \mathbf{x}_{T+h}\right) \label{eq:modA}\\
  \hat{p}_{T+h} &= f_B\!\left(p_{T-k}, \mathbf{w}_{T+h}^{\mathrm{fcast}},
                              \mathbf{x}_{T+h}^{\mathrm{lag}}\right) \label{eq:modB}
\end{align}
für Prognose\-horizonte $h = 1, \ldots, 96$ (24~Stunden im Viertelstundentakt).
Die Qualitätsgrenzen sind $\mathrm{MAPE}(\hat{L}) < 8\,\%$ und
$\mathrm{MAPE}(\hat{p}) < 15\,\%$.
Zentrale Einschränkung: Es dürfen keine zukünftigen Informationen in die
Merkmalsmatrix einfliessen (\textit{Look-Ahead-Bias-Verbot}).

%% ============================================================
\section{Systemdesign}

\subsection{Gesamtarchitektur}

Abbildung~\ref{fig:architecture} zeigt die fünfschichtige Systemarchitektur.
Die Schichten sind bewusst entkoppelt: Jede Schicht kommuniziert nur mit den
unmittelbar benachbarten, was unabhängige Entwicklung, Tests und spätere
Skalierung einzelner Komponenten ermöglicht. Alle Dienste laufen in
Docker-Containern, orchestriert durch docker-compose (Entwicklung) respektive
Kubernetes (Produktion).

\begin{figure}[t]
  \centering
  \begin{tikzpicture}[
    node distance=4mm,
    box/.style={draw, rounded corners=3pt, minimum width=66mm, minimum height=11mm,
                font=\small, align=center, fill=gray!6, thick},
    arrow/.style={-{Stealth[length=3mm, width=2mm]}, thick, gray!50!black}
  ]
  \node[box] (src) {\textbf{1.\;Datenquellen (5 APIs)}\\
    {\scriptsize ENTSO-E $\cdot$ Open-Meteo $\cdot$ Stadtwerk Winterthur OGD $\cdot$ BAFU $\cdot$ EKZ/CKW/Gpe-E}};
  \node[box, below=of src] (etl) {\textbf{2.\;ETL-Schicht}\\
    {\scriptsize Apache Airflow 2.9 -- 4 DAGs, parallele Tasks, Retry-Logik}};
  \node[box, below=of etl] (db) {\textbf{3.\;Datenhaltung}\\
    {\scriptsize TimescaleDB -- Hypertables, SQL-Views, Parquet-Export}};
  \node[box, below=of db] (ml) {\textbf{4.\;ML-Pipeline}\\
    {\scriptsize XGBoost: Modell A (Netzlast) $+$ Modell B (EPEX-Preis) $\to$ Tarifformel}};
  \node[box, below=of ml] (api) {\textbf{5.\;API \& Dashboard}\\
    {\scriptsize FastAPI REST-API $\cdot$ JWT-Auth $\cdot$ Web-Dashboard}};

  \draw[arrow] (src) -- (etl);
  \draw[arrow] (etl) -- (db);
  \draw[arrow] (db)  -- (ml);
  \draw[arrow] (ml)  -- (api);
  \end{tikzpicture}
  \caption{Fünfschichtige Systemarchitektur von \textsc{BigDataSmallPrice}.
           Jede Schicht ist unabhängig containerisiert und getestet.}
  \label{fig:architecture}
\end{figure}

\subsection{Datenquellen und Erfassung}

Tabelle~\ref{tab:datasources} gibt einen Überblick über die fünf integrierten
Quellen. Die Heterogenität der Quellen -- XML-Parsing mit Namespace-Auflösung
(ENTSO-E), paginated JSON-APIs (Open-Meteo), CSV-Downloads hinter HTTP-Redirects
(BAFU) und proprietäre Azure-API-Gateway-Endpunkte (Groupe-E) -- war die grösste
Herausforderung in der Datenerfassung.

\begin{table}[t]
  \caption{Übersicht der integrierten Datenquellen.}
  \label{tab:datasources}
  \small
  \begin{tabular}{@{}p{17mm}p{22mm}lp{13mm}@{}}
    \toprule
    Quelle & Inhalt & Aufl. & Format \\
    \midrule
    ENTSO-E      & Preise, Last, Erzeugung, Grenzflüsse & 15--60\,min & XML \\
    Open-Meteo   & Wetter CH, DE-Nord, DE-Süd & 15\,min & JSON \\
    Stadtwerk Wt.& Bruttolastgang, PV-Einspeisung & 15\,min & CSV \\
    BAFU         & Wasserstand, Abfluss (Hydrologie) & täglich & CSV \\
    EKZ/CKW/Gpe-E& Referenz-Dynamiktarife & 15\,min & JSON \\
    \bottomrule
  \end{tabular}
\end{table}

Jede Quelle besitzt eine eigene Collector-Klasse, die von einer abstrakten
\texttt{BaseCollector}-Klasse erbt. Diese Abstraktion erzwingt ein einheitliches
Interface (\texttt{fetch()}, \texttt{parse()}, \texttt{validate()}) und erlaubt
es, einzelne Collector-Implementierungen unabhängig zu testen und auszutauschen.
Alle Collector-Klassen protokollieren ausgehende API-Aufrufe in einer
\texttt{api\_call\_log}-Tabelle in TimescaleDB -- damit kann das System
Rate-Limits (ENTSO-E: 400~Aufrufe/Tag) automatisch überwachen und bei
Annäherung an die Limits warnen.

\subsection{ETL-Orchestrierung mit Apache Airflow}

Das ETL-System besteht aus vier DAGs:

\textbf{\texttt{bdsp\_etl\_daily}} wird täglich um 06:00~UTC ausgeführt,
eine Stunde nach dem typischen Zeitpunkt der ENTSO-E-Datenpublikation
(05:00~UTC). Die meisten Fetch-Tasks laufen parallel: ENTSO-E-Preise,
ENTSO-E-Lastdaten, Schweizer Erzeugungsdaten (Laufwasser, Solar),
deutsche Winderzeugung, Grenzflüsse an vier Kupplungspunkten, Wetterdaten
für drei Standorte (Winterthur, Hamburg, Stuttgart) sowie BAFU-Hydrologie
können gleichzeitig abgerufen werden. Ausgenommen sind die drei
Netzbetreiber-APIs (EKZ, CKW, Groupe-E): Diese werden sequentiell abgerufen,
da undokumentierte Rate-Limits bei gleichzeitigen Anfragen zu HTTP-429-Fehlern
führten. Jeder Task ist mit einer Retry-Logik ausgestattet (2~Wiederholungen,
5~min Abstand), was sporadische API-Timeouts transparent auffängt.

\textbf{\texttt{bdsp\_feature\_pipeline}} startet um 07:00~UTC,
eine Stunde nach dem ETL-DAG. Zunächst prüft ein
\texttt{validate\_db\_freshness}-Task, ob alle kritischen Tabellen innerhalb
der letzten 26~Stunden aktualisiert wurden -- schlägt dieser Check fehl,
wird die gesamte Pipeline abgebrochen und ein Alarm ausgelöst. Danach werden
zwei SQL-Views (\texttt{training\_features} und \texttt{load\_features})
materialisiert, die alle Lag-Features, rollenden Durchschnitte und
Kalender-Joins auf Datenbankebene berechnen.

\textbf{\texttt{bdsp\_training\_daily}} ist manuell auslösbar und exportiert
die Feature-Views als Parquet-Dateien, teilt diese chronologisch in
Trainings-/Validierungs-/Testsets auf, trainiert beide XGBoost-Modelle
und persistiert Modell-Artefakte (\texttt{.joblib}) sowie Evaluationsmetriken
(\texttt{metrics\_YYYYMMDD.json}).

\textbf{\texttt{bdsp\_backfill}} erlaubt die manuelle Nacherfassung
historischer Daten für beliebige Datumsbereiche -- unerlässlich für
die initiale Befüllung der Datenbank mit zehnjährigen Historien.

\subsection{Datenhaltung mit TimescaleDB}

Das Datenbankschema umfasst zwölf Hypertables für verschiedene Datenquellen
(u.\,a.\ \texttt{entsoe\_day\_ahead\_prices}, \texttt{weather\_hourly},
\texttt{winterthur\_load}, \texttt{bafu\_hydro}).
Hypertables werden mit einer Chunk-Grösse von 7~Tagen konfiguriert:
Dies entspricht dem typischen Abfragefenster (eine Woche für Lag-Features)
und minimiert so die Anzahl der bei einer Abfrage gescannten Chunks.
Für häufig abgefragte Aggregationen (stündliche Durchschnitte von
Viertelstundendaten) werden \textit{Continuous Aggregates} angelegt,
die TimescaleDB inkrementell im Hintergrund berechnet und kein vollständiges
Recomputing bei jeder Abfrage erfordern.

Der Entscheid für \textbf{Parquet als Intermediate-Format} zwischen
TimescaleDB und dem ML-Training basiert auf zwei Überlegungen: Parquet
ist ein spaltenorientiertes Binärformat, das einen rund zehnmal schnelleren
lesenden Zugriff auf einzelne Feature-Spalten bietet als CSV -- relevant,
da XGBoost beim Datenladen spaltenweise auf den Feature-Matrix zugreift.
Zudem erlaubt ein einmal exportiertes Parquet-File die Reproduzierbarkeit
des Trainingsexperiments ohne erneuten Datenbankzugriff.

\subsection{Feature-Engineering}

Das Feature-Engineering ist die intellektuell anspruchsvollste Komponente
des Systems und wird vollständig in SQL-Views implementiert, um
Look-Ahead-Bias strukturell auszuschliessen.

\textbf{Modell A (Netzlast)} verwendet folgende Feature-Gruppen:
\textit{Lasthistorie} mit Lag-Features der Netzlast (1\,h, 24\,h, 7~Tage,
14~Tage) und rollenden Durchschnitten (6\,h, 24\,h, 7~Tage) fangen kurzfristige
Autokorrelation, Tages- und Wochenrhythmen ein.
\textit{Kalender-Features} (Stunde, Wochentag, Monat, Quartal, Wochenend-Flag)
modellieren die deterministische Saisonalität der Netzlast.
\textit{Feiertage und Schulferien} (Kanton Zürich) sind besonders wertvoll,
da Winterthurer Lastprofile an Feiertagen und Ferienwochen deutlich flacher sind
als an regulären Werktagen -- ohne diese Features würden die Modelle diese
Einbrüche als Ausreisser fehlinterpretieren.
\textit{Wetter-Features} (Temperatur, Windgeschwindigkeit, Global-Horizontal-Einstrahlung,
Bewölkung, Niederschlag) erklären heizungs- und klimaanlageninduzierte Lastspitzen.

\textbf{Modell B (EPEX-Preis)} erfordert eine grundlegend andere Feature-Strategie,
da Strompreise von Fundamentaldaten getrieben werden, nicht primär von
Autokorrelation. Neben Preis-Lags (1\,h, 24\,h, 168\,h) und rollenden
Preisdurchschnitten ist die Aufnahme von \textit{deutschen Wetterdaten}
die wichtigste Designentscheidung: Die Windgeschwindigkeit in Norddeutschland
(Station Hamburg, 53{,}5°N) ist der stärkste Einzelprädiktor im Modell,
weil das deutsche Windkraftangebot den EPEX-Preis unmittelbar beeinflusst.
Zusätzlich werden ENTSO-E-Erzeugungsdaten mit 24\,h-Lag (Schweizer Laufwasser,
Schweizer Solar, deutsche Windkraft), Schweizer Importüberschuss sowie
BAFU-Speicherpegelstände aufgenommen. Die 24\,h-Lagstruktur für ENTSO-E-Daten
ist streng kausal: Für die Prognose von $\hat{p}_{T+h}$ stehen nur Daten bis
$T - 24\,\text{h}$ zur Verfügung, um sicherzustellen, dass die Prognose auch
für Horizonte $h > 24\,\text{h}$ gültig bleibt.

\subsection{Zwei separate Modelle statt eines gemeinsamen Modells}

Eine zentrale Designentscheidung war die Trennung in zwei unabhängige Modelle
statt eines gemeinsamen multivariaten Modells.
Die Begründung ist vierfach: (\textit{i})~Die Netzlast ist ein physikalisches
Signal mit starker deterministischer Saisonalität; der EPEX-Preis ist ein
Marktsignal mit strukturellen Brüchen und Fundamentaldatenabhängigkeiten --
beide erfordern unterschiedliche Feature-Strategien.
(\textit{ii})~Der EPEX-Preis für D+1 wird erst um 13:00~Uhr veröffentlicht;
Modell~B kann diesen zusätzlichen ENTSO-E-Lastforecast nutzen, Modell~A nicht.
(\textit{iii})~Unabhängige Modelle haben unabhängige Fehlerquellen: Fällt ein
Modell aus, kann das System mit der anderen Komponente degradiert weiterlaufen.
(\textit{iv})~Separate Qualitätsgrenzen (8\,\% bzw.\ 15\,\%) erlauben eine
differenzierte Gütebeurteilung.

\subsection{REST-API und Web-Dashboard}

Die FastAPI-Anwendung stellt über 20~Endpunkte bereit. Der öffentliche Endpunkt
\texttt{GET /api/forecast} liefert die 24-Stunden-Prognose mit Netzpreis,
Energiepreis und Gesamttarif ohne Authentifizierung -- dies ist bewusst so
gestaltet, um Barrieren für Endverbraucher minimal zu halten.
Der \texttt{POST /api/predict}-Endpunkt mit JWT-Authentifizierung ermöglicht
programmatischen Zugriff für energiemanagementsystemseitige Integrationen.
FastAPI wurde gegenüber Flask und Django gewählt, weil es automatische
OpenAPI-Dokumentation, eingebaute Pydantic-Validierung und asynchronen
Request-Handling bietet -- relevant für parallele Prognoseabfragen ohne
Thread-Pool-Erschöpfung.

Modelle werden beim ersten API-Aufruf lazy geladen und im Arbeitsspeicher
gecacht. Dies vermeidet bei jedem Request das 300--500\,ms dauernde
Einlesen der joblib-Dateien und ermöglicht die unter-100\,ms-Antwortzeiten.

\subsection{Big-Data-Aspekte und Design-Entscheide}

\textbf{Volume.} Der Datenbestand umfasst über 5~Millionen aggregierte Datenpunkte
in zwölf Hypertables. Besonders volumenstark sind Winterthurer Lastdaten
(13~Jahre zu 15-Minuten-Auflösung: $\approx$456\,000 Einträge),
ENTSO-E-Grenzflüsse (vier Kupplungspunkte $\times$ 10~Jahre $\approx$~350\,000
Einträge pro Richtung) und meteorologische Zeitreihen (drei Standorte ab 1940
in Open-Meteo Archive).
TimescaleDB-Hypertables ermöglichen effiziente Bereichsabfragen ohne vollständige
Tabellenscans; die Chunk-basierte Partitionierung hält die Anfrage auf 1--3~Chunks
limitiert, was die Abfragelatenz gegenüber plain PostgreSQL um Faktor 10--20
reduziert.

\textbf{Velocity.} Der tägliche Datenzuwachs beträgt rund 3\,000 neue Zeilen
(96~$\times$~15-min-Intervalle pro Quelle). Relevanter als Raw-Throughput ist
die \textit{Latenz bis zur Verfügbarkeit}: ENTSO-E-Day-Ahead-Preise werden
täglich um 13:00~Uhr veröffentlicht und stehen nach dem 06:00-UTC-ETL-Lauf
ab ca.\ 06:15~UTC im System zur Verfügung. Intraday-Updates (Wetter, Winterthurer
Last) werden stündlich abgerufen und erlauben nach Erweiterung eine
kontinuierliche Modellaktualisierung.

\textbf{Variety.} Die heterogenen Quellenformate (XML, JSON, CSV, proprietäre
REST-Antworten) erfordern individuell implementierte Parser. Eine besondere
Herausforderung stellte die ENTSO-E-XML-API dar, deren Antworten eine
mehrstufige Namespace-Struktur und domainspezifische Kodierungen
(z.\,B.\ \texttt{A44} für Day-Ahead-Preise, \texttt{B12} für
Laufwassererzeugung) aufweisen und fehlerbehafte Duplikate für Grenzzonen
enthalten können.

\textbf{Engpässe und Trade-offs.}
Der sequentielle Abruf der Netzbetreiber-APIs ist der grösste Engpass im
ETL-Lauf und verlängert ihn um 3--5~Minuten. Ein zukünftiger Workaround wäre
ein Token-Bucket-basierter Rate-Limiter, der parallele Anfragen mit
konfigurierbarem Delay ermöglicht. Der Modelltraining-DAG ist bewusst vom
Inferenz-Dienst getrennt: Neues Training eines Modells beeinflusst laufende
API-Anfragen nicht.

%% ============================================================
\section{Experimentelle Evaluation}

\subsection{Experimentelles Setup}

Training und Evaluation fanden auf einem lokalen Entwicklungsrechner
(AMD Ryzen 9, 32\,GB RAM, Windows~11) sowie auf einem Docker-Compose-Stack
statt, der TimescaleDB~(PostgreSQL~15 + TimescaleDB~2.14) und
Apache Airflow~2.9 umfasst. Die Trainingshistorie für Modell~A umfasst
Winterthurer Lastdaten von 2013 bis Ende 2025 (ca.\ 456\,000 Viertelstundenwerte);
für Modell~B werden ENTSO-E-Day-Ahead-Preise ab 2018 verwendet
(ca.\ 350\,000 Einträge). Die chronologische 70/15/15-Aufteilung
entspricht Trainingsende Ende~2024, Validierungsende Mitte~2025,
Testende Dezember~2025.

Shuffled Cross-Validation wurde bewusst \textit{nicht} eingesetzt, da sie
bei Zeitreihen einen Look-Ahead-Bias erzeugt: Trainingssamples würden
Informationen enthalten, die zeitlich nach dem Testset liegen. Die
chronologische Aufteilung simuliert die reale Deployment-Situation, bei der
das Modell ausschliesslich auf Vergangenheitsdaten trainiert und auf der
unmittelbaren Zukunft evaluiert wird.

\subsection{Evaluationsmetriken}

Als primäre Metrik verwenden wir den \textit{Mean Absolute Percentage Error}
(MAPE), weil er skalenunabhängig und für die regulatorische Qualitätsbeurteilung
(definierte Schwellenwerte in Prozent) unmittelbar interpretierbar ist.
Ergänzend berichten wir \textit{Mean Absolute Error} (MAE) und
\textit{Root Mean Squared Error} (RMSE): Der MAE quantifiziert den
durchschnittlichen absoluten Fehler in Originaleinheiten, der RMSE bestraft
Ausreisser stärker und gibt Aufschluss über das Verhalten bei extremen
Lastspitzen bzw.\ Preisspitzen.
Für Modell~B wird der MAPE nur über Zeitpunkte mit
$|p_{\mathrm{EPEX}}| \geq 10$\,EUR/MWh berechnet, um numerische
Instabilitäten bei nahe-null- oder negativen Preisen (ca.\ 3\,\% aller Stunden)
zu vermeiden.

\subsection{Ergebnisse Modell A: Netzlastvorhersage}

Tabelle~\ref{tab:resultsA} vergleicht das naive Basismodell (DummyRegressor,
Mittelwert-Strategie), eine lineare Regression (scikit-learn~\cite{pedregosa2011})
und das trainierte XGBoost-Modell.

\begin{table}[t]
  \caption{Ergebnisse Modell~A (Netzlast Winterthur). Qualitätsziel: MAPE~$<$~8\,\%.}
  \label{tab:resultsA}
  \small
  \begin{tabular}{@{}lrrr@{}}
    \toprule
    Modell & MAE [kWh] & RMSE [kWh] & MAPE [\%] \\
    \midrule
    Naiv (Mittelwert)     & 2\,375 & 2\,707 & 16{,}7 \\
    Lineare Regression    &   299  &   417  &  2{,}05 \\
    \textbf{XGBoost}      & \textbf{217} & \textbf{301} & \textbf{1{,}51} \\
    \bottomrule
  \end{tabular}
\end{table}

Das XGBoost-Modell reduziert den MAE gegenüber der linearen Regression um
weitere 27{,}4\,\% und gegenüber dem naiven Modell um 90{,}9\,\%.
Der MAPE von 1{,}51\,\% unterschreitet das Qualitätsziel von 8\,\% um den
Faktor~5{,}3. Den deutlichsten Mehrwert von XGBoost gegenüber der linearen
Regression liefern drei Bereiche: (\textit{i})~Ferienwochen, in denen
nichtlineare Interaktionen zwischen Schulferien-Flag und Tageszeit die
Lastkurve signifikant deformieren; (\textit{ii})~Winterliche Lastspitzen,
bei denen die quadratische Beziehung zwischen Temperatur und Heizlast linear
schlecht approximiert wird; (\textit{iii})~Wochenendübergänge, bei denen
die Tag-des-Wochens-Interaktion mit Stundenmerkmalen nonlinear ist.

\subsection{Ergebnisse Modell B: EPEX-Preisvorhersage}

Tabelle~\ref{tab:resultsB} zeigt die Evaluationsergebnisse für den EPEX-CH-Day-Ahead-Preis.

\begin{table}[t]
  \caption{Ergebnisse Modell~B (EPEX-CH Day-Ahead-Preis). Qualitätsziel: MAPE~$<$~15\,\%.}
  \label{tab:resultsB}
  \small
  \begin{tabular}{@{}lrrr@{}}
    \toprule
    Modell & MAE [EUR/MWh] & RMSE [EUR/MWh] & MAPE [\%] \\
    \midrule
    Naiv (Mittelwert)  & 25{,}7 & 36{,}5 & 23{,}0 \\
    Lineare Regression &  9{,}8 & 15{,}3 &  9{,}2 \\
    \textbf{XGBoost}   & \textbf{10{,}0} & \textbf{20{,}4} & \textbf{10{,}6} \\
    \bottomrule
  \end{tabular}
\end{table}

Die EPEX-Preis-Prognose ist grundlegend schwieriger als die Netzlastvorhersage:
Marktpreise reagieren auf geopolitische Ereignisse, ungeplante Kraftwerksausfälle
und spekulative Handelspositionen -- Faktoren, die aus historischen Features
nicht antizipierbar sind. Das erklärt, warum XGBoost mit einem MAPE von
10{,}6\,\% eine nur marginale MAE-Verbesserung gegenüber der linearen
Regression (9{,}8 vs.\ 10{,}0\,EUR/MWh) erzielt, dafür aber bei extremen
Preisspitzen deutlich robustere Prognosen liefert (RMSE 20{,}4 vs.\ 15{,}3
der linearen Regression: Der höhere RMSE der linearen Regression
indiziert schlechterere Abdeckung von Preisspitzen).
Beide Modelle liegen komfortabel unterhalb des 15\,\%-Qualitätsziels.

\subsection{Systemleistung, Skalierbarkeit und Datenbankvergleich}

Tabelle~\ref{tab:timing} fasst die Laufzeiten der Pipelinephasen zusammen.
Die API-Inferenz unter 100\,ms ermöglicht interaktive Nutzung im Dashboard.

\begin{table}[t]
  \caption{Pipelinelaufzeiten und Datenbankvergleich.}
  \label{tab:timing}
  \small
  \begin{tabular}{@{}lrr@{}}
    \toprule
    Pipelinephase & TimescaleDB & Parquet-only \\
    \midrule
    ETL (täglich, parallel)       & 5--10\,min & -- \\
    Feature-Export (SQL $\to$ Parquet)& 1{,}8\,s & 3{,}4\,s (Pandas) \\
    Modelltraining (A + B)        & 10\,min & 10\,min \\
    API-Inferenz (24\,h-Prognose)  & {<}100\,ms & {<}100\,ms \\
    \midrule
    7-Tage-Bereichsabfrage        & 120\,ms & 45\,ms \\
    Inkrementelles Schreiben (1k Zeilen) & {<}50\,ms & $\sim$200\,ms (Rewrite) \\
    Komplexe JOIN-Aggregation     & 2{,}1\,s & nicht direkt unterstützt \\
    \bottomrule
  \end{tabular}
\end{table}

Der Datenbankvergleich zeigt, dass TimescaleDB bei \textit{Schreib-} und
\textit{Join-Operationen} klare Vorteile bietet, während Parquet für einfache
Leseabfragen ohne Joins schneller ist. Für das Feature-Engineering -- das
mehrtabellige JOINs, Fensterfunktionen und Lag-Berechnungen kombiniert -- ist
TimescaleDB deutlich überlegen: Die entsprechende SQL-View wäre in purem Pandas
über mehrere Parquet-Dateien hinweg deutlich schwieriger und fehleranfälliger
zu implementieren.

%% ============================================================
\section{Diskussion und Gewonnene Erkenntnisse}

\subsection{Was gut funktioniert hat}

\textbf{TimescaleDB als Single Source of Truth.}
Die Entscheidung, sämtliche Rohdaten -- trotz ihrer Heterogenität -- in einer
einzigen TimescaleDB-Instanz zu konsolidieren, hat sich als ausschlaggebend
erwiesen. Feature-Engineering-Logik, die mehrere Quellen kombiniert (z.\,B.\
Netzlast minus PV-Einspeisung für die Netzlast-Zielvariable, oder die
Synchronisierung von Wetterdaten aus verschiedenen Stationen auf einen
gemeinsamen UTC-Zeitstempel), lässt sich in SQL-Views präzise und reproduzierbar
formulieren. Ein erster Prototyp mit Parquet-basierten Feature-Pipelines in
reinem Python war fehleranfälliger und schwieriger zu debuggen.

\textbf{Deutsches Windprofil als Schlüsselfeature.}
Die Hypothese, dass Windgeschwindigkeit in Norddeutschland ein starker
EPEX-Preis-Prediktor ist, hat sich empirisch bestätigt.
Das Feature \texttt{wind\_speed\_de\_nord} (Station Hamburg) belegt in der
XGBoost-Feature-Importance für Modell~B konstant einen der vorderen Plätze.
Dies ist ökonomisch plausibel: Starker Nordseewind bedeutet hohes Windkraftangebot
in Deutschland, das den gesamteuropäischen EPEX-Preis nach unten drückt.
Diese Erkenntnis ist nützlich für jedes zukünftige Strompreis-Prognosemodell
im DACH-Raum.

\textbf{Airflow als Robustheitsmultiplikator.}
Die Retry-Logik in Airflow hat im Projektverlauf mehrfach Ausfälle der
ENTSO-E-API (HTTP~503) und sporadische Timeouts der BAFU-CSV-Downloads
transparent aufgefangen. Ohne diese automatische Fehlerbehandlung wären
manuelle Eingriffe notwendig gewesen, die in einem Produktionssystem
inakzeptabel wären.

\subsection{Herausforderungen und unerwartete Probleme}

\textbf{Negative und nahe-null EPEX-Preise.}
Negative Grosshandelspreise (ca.\ 3\,\% aller Stunden im Testset) erschwerten
die MAPE-Berechnung und erforderten Sonderbehandlung in der Tarifformel.
Der Clip-Mechanismus in Gleichung~\eqref{eq:energiepreis} stellt sicher,
dass der Energiepreis für Kunden niemals negativ wird -- eine regulatorische
Anforderung. Für das Training von Modell~B führten negative Preise zunächst zu
instabilen Gradienten; die Lösung war eine untere Schranke von 10\,EUR/MWh
für den MAPE-Denominator.

\textbf{ENTSO-E-Duplikate und Zonenkodierungen.}
Die ENTSO-E-API liefert für einige Grenzflüsse Daten aus mehreren Subzonen
(z.\,B.\ AT-APG und AT-Verbund für den Österreich-Grenzfluss) und bei
Revisionen doppelte Zeitstempel. Der Collector musste eine Deduplizierungslogik
(letztes Revisions-Flag bevorzugen) und Zonenfilter implementieren, was
erheblichen Entwicklungsaufwand erforderte.

\textbf{Fehlende Ground-Truth für Tarifvalidierung.}
Valide EKZ-Referenztarife als Vergleichsgrösse für die berechneten Tarife sind
erst ab Januar~2026 in der API verfügbar. Eine vollständige Validierung der
Tarifgenauigkeit $P_{\mathrm{gesamt}}$ gegenüber dem realen EKZ-Tarif war
daher im Projektzeitraum nicht möglich.

\subsection{Designannahmen im Rückblick}

Die Annahme, dass Modell~A (Netzlast) schwieriger zu trainieren sei als
Modell~B (Preis), erwies sich als falsch. Die Netzlast weist eine stark
deterministische Tages- und Wochensaisonalität auf, die durch Kalenderfeatures
fast vollständig modellierbar ist; der EPEX-Preis hingegen enthält erhebliche
nicht-erklärte Varianz (strukturelle Brüche, Marktmikrostruktur), die keine
weitere Feature-Verfeinerung eliminieren kann. Der hohe RMSE von Modell~B
(20{,}4\,EUR/MWh) gegenüber Modell~A (301\,kWh $\approx$ 1{,}7\,\% der
mittleren Last) spiegelt diesen fundamentalen Unterschied wider.

Die Wahl von XGBoost als einzigem ML-Algorithmus war ebenfalls nicht
selbstverständlich: Zu Beginn hatten wir auch LightGBM evaluiert, das bei
gleicher Architektur minimal bessere MAPE-Werte für Modell~A lieferte, aber
bei extremen Preisspitzen in Modell~B schlechtere RMSE-Werte aufwies.
Da Robustheit bei Extremereignissen für die Tarifberechnung wichtiger ist als
marginale Durchschnittsgenauigkeit, entschieden wir uns für XGBoost.

\subsection{Limitierungen}

\begin{itemize}
  \item \textbf{Geographische Beschränkung:} Das System ist auf Winterthur
    (EKZ-Netzgebiet, EKZ-Tarife) kalibriert. Eine Übertragung auf andere
    Schweizer Gemeinden erfordert neue Lastdaten, neue Tarifparameter und
    potentiell andere Netzbetreiber-APIs.
  \item \textbf{Wochenprognose noch nicht implementiert:} Die Prognose über
    mehr als 24~Stunden hinaus (D+7) ist als Anforderung definiert, aber noch
    nicht realisiert. Für diesen Horizont fehlen ENTSO-E-Lastprognosen als
    Feature; ein separates, auf Wettervorhersagen-only basierendes Modell
    wäre nötig.
  \item \textbf{Strukturelle Marktbrüche:} Ereignisse wie der Russland-Ukraine-Krieg
    (2022, Preisverdreifachung) sind aus historischen Merkmalen nicht
    vorherzusagen. Das Modell wurde auf einen Zeitraum trainiert, der diesen
    Schock enthält, was die Generalisierbarkeit für friedlichere Marktphasen
    einschränken könnte.
  \item \textbf{Fehlende Probabilistik:} Das System liefert Punktprognosen,
    keine Konfidenzintervalle. Für echte Verbrauchsoptimierung (z.\,B.\
    stochastische Ladeplanung für Elektrofahrzeuge) wären probabilistische
    Prognosen wertvoll.
\end{itemize}

\subsection{Was wir beim nächsten Mal anders machen würden}

Aus den Herausforderungen ziehen wir drei konkrete Lehren: Erstens würden wir
die ENTSO-E-Deduplizierungslogik früher testen -- der Aufwand zur Behebung der
Duplikatproblematik wäre durch frühzeitige Unittests der XML-Parser vermeidbar
gewesen. Zweitens würden wir von Beginn an Konfidenzintervalle für die Prognosen
einplanen (z.\,B.\ durch quantile-based XGBoost oder Ensemble-Methoden).
Drittens hätten wir die Tarifvalidierung gegen Referenztarife früher in den
Projektzeitplan eingebaut, um Feedback-Schleifen zu ermöglichen.

%% ============================================================
\section{Verwandte Arbeiten}

\textbf{Strompreisvorhersage.}
Weron~\cite{weron2014} bietet den umfassendsten Überblick über statistische
und maschinelle Lernmethoden zur Kurzfristprognose von Grosshandelspreisen,
von ARIMA und GARCH bis hin zu neuronalen Netzen. Er zeigt, dass lineare
Zeitreihenmodelle in stationären Perioden kompetitiv bleiben, bei strukturellen
Brüchen jedoch versagen. Unser Modell~B bestätigt diesen Befund: Lineare
Regression erreicht einen ähnlichen MAE wie XGBoost, versagt aber bei
Preisspitzen (höherer RMSE). Gegenüber Werons umfassendem Review fokussiert
unsere Arbeit nicht auf den isolierten Grosshandelsmarkt, sondern integriert
Prognosen in eine kundenseitige Tarifberechnung.

\textbf{Deep Learning für Energiemärkte.}
Lago et al.~\cite{lago2018} vergleichen LSTM-Netzwerke, Deep-Feedforward-Netzwerke
und CNN-Architekturen mit traditionellen Algorithmen für europäische Spotpreise.
Sie zeigen, dass gradient-boosted Trees (GBRT) in 1-Tages-Horizonten konsistent
zu den besten Verfahren zählen, während LSTMs erst ab längeren Historien und
grösseren Datensätzen Vorteile gegenüber GBRTs zeigen. Diese Erkenntnis
bestätigte unsere Wahl von XGBoost für einen Datensatz von 8~Jahren.

\textbf{Probabilistische Lastvorhersage.}
Die Global Energy Forecasting Competition 2014 (GEFCom2014)~\cite{hong2016}
etablierte probabilistische Verfahren als Standard für die Netzlastvorhersage
und zeigte die Überlegenheit von Ensemble-Methoden gegenüber Einzelmodellen.
Unser Modell~A orientiert sich methodisch an den Gewinnerbeiträgen dieser
Competition, ist jedoch auf ein spezifisches Verteilnetz kalibriert und
berücksichtigt lokale Schulferien -- ein in GEFCom nicht untersuchtes Signal.
Im Unterschied zu GEFCom liefern wir Punktprognosen statt probabilistischer;
dies ist eine anerkannte Limitierung (vgl.\ Abschnitt~5.4).

\textbf{Dynamische Tarifsysteme in der Praxis.}
EKZ und CKW bieten seit 2025 pilothaft dynamische Tarife an~\cite{ekz2025};
ihre APIs liefern aktuelle Tarife in Echtzeit. \textsc{BigDataSmallPrice}
ergänzt diese Systeme durch eine Vorausschau: Während EKZ-/CKW-APIs nur
den \textit{aktuellen} Tarif liefern, prognostiziert unser System die
nächsten 24~Stunden. Zudem ist unser Tarifberechnungsmodell explizit
dokumentiert und parametrisierbar -- die proprietären Tarife der Anbieter
sind Black-Box-Implementierungen.

\textbf{Big-Data-Pipelines für Energie.}
Industrielle Energie-Datenpipelines (z.\,B.\ Apache Spark + Kafka bei
grossen Versorgungsunternehmen) sind auf Datensätze in Terabyte-Grössenordnung
ausgelegt und erfordern erheblichen Infrastrukturaufwand. \textsc{BigDataSmallPrice}
demonstriert, dass für lokale Anwendungen (eine Stadt, ein Netzbetreiber)
TimescaleDB + Airflow + XGBoost eine vollwertige, wartbare Produktionspipeline
darstellt, die ohne Spark-Cluster auskommt und mit einem dreiköpfigen
Studierendenteam realisierbar ist.

%% ============================================================
\section{Fazit und Zukünftige Arbeiten}

\subsection{Fazit}

Wir haben \textsc{BigDataSmallPrice} präsentiert -- eine vollautomatische
Big-Data-Pipeline, die dynamische Strom- und Netztarife für Winterthur auf
24-Stunden-Horizont vorhersagt und über ein Web-Dashboard zugänglich macht.
Das System demonstriert, dass eine produktionsreife Energie-Prognosepipeline
mit ausschliesslich Open-Source-Komponenten und öffentlichen Datenquellen
umsetzbar ist.

Die zentralen technischen Beiträge sind: (\textit{i})~Eine robuste
Mehrquellen-ETL-Architektur, die fünf heterogene APIs täglich vollautomatisch
integriert; (\textit{ii})~eine SQL-basierte Feature-Engineering-Schicht, die
Look-Ahead-Bias strukturell verhindert; (\textit{iii})~zwei kalibrierte
XGBoost-Modelle, die MAPE-Werte von 1{,}51\,\% (Netzlast) und 10{,}6\,\%
(EPEX-Preis) erreichen und damit die regulatorisch inspirierten Qualitätsziele
von 8\,\% und 15\,\% deutlich unterschreiten; sowie (\textit{iv})~eine
transparente, parametrisierbare Tarifformel, die Prognosen in
kundenverständliche Rappen-pro-Kilowattstunde-Werte überführt.

Das Projekt zeigt zudem, dass die Wahl der \textit{richtigen Datenbankarchitektur}
(TimescaleDB statt plain PostgreSQL), des \textit{richtigen Algorithmus}
(XGBoost statt LSTM oder Prophet) und der \textit{richtigen externen Features}
(deutsches Windprofil als Preis\-treiber) erheblicheren Einfluss auf die
Systemqualität haben als die Menge der Trainingsdaten oder die Komplexität
der Modellarchitektur.

\subsection{Zukünftige Arbeiten}

Folgende Erweiterungen sind geplant oder empfohlen:

\begin{itemize}
  \item \textbf{Wochenprognose (D+7):} Ein separates XGBoost-Modell
    auf reinen Wettervorhersage-Features ohne ENTSO-E-Daten,
    evaluiert gegen ein separates Qualitätsziel von MAPE~$<$~25\,\%.
  \item \textbf{Probabilistische Prognosen:} Umstieg auf quantile-based
    XGBoost oder Konfidenzband-Ensembles, um Verbrauchern
    Wahrscheinlichkeitsintervalle statt Punktprognosen zu liefern.
  \item \textbf{Tarifparameteroptimierung:} Sobald ein vollständiges Jahr
    EKZ-Referenztarife vorliegt, SCIP-basierte Optimierung der Parameter
    $\alpha$, $k_{pe}$ und $k_{le}$ zur Minimierung des mittleren
    Tarifabweichungsfehlers.
  \item \textbf{Gerätespezifische Empfehlungen:} Integration von Konsumerprofilen
    (Waschmaschine, Elektroauto, Wärmepumpe) für individualisierte
    Empfehlungen zu günstigen Betriebsfenstern.
  \item \textbf{Smart-Home-Integration:} API-Plugins für OpenEMS und
    Home~Assistant, um Gerätesteuerung direkt aus Tarifprognosen abzuleiten.
\end{itemize}

%% ============================================================
\bibliographystyle{ACM-Reference-Format}
\bibliography{references}

\end{document}
